\chapter{实验与分析}

\section{实验准备}

\subsection{环境配置}
本文实验环境为：处理器为 Intel Core i7-12700h、运行内存为
16G、GPU 为 NVIDlA GeForce RTX 3060 Laptop GPU。本文实验基于 PyCharm 和 vscode平台对模型算法进行训练和测试，操作系统为 Windows11，配套环境为 python3.8、CUDA11.1、carla 0.9.15，仿真环境为 CarlaUE4。
\subsection{实验数据集与参数设置}

本次实验基于CARLA仿真平台搭建测试场景，分别启用CARLA原生物理引擎（基于Bullet游戏引擎）与Chrono多体物理引擎，保证两个测试场景的参数完全一致，包括车辆模型、初始状态、行驶路径、环境参数及仿真时长等。实验仿真时长设定为20秒，其中0-15秒为车辆正常行驶阶段，包含加速、减速、转向等常规操作，15-20秒为车辆碰撞测试阶段，模拟极端工况下的车辆动力学响应。实验过程中，同步采集两个引擎输出的IMU核心数据，包括平动速度（X/Y/Z轴）、平动加速度（X/Y/Z轴）及角速度（X/Y/Z轴），形成对比数据集用于后续分析。
\section{评价指标}
本文以“相同行动路径”为核心前提（即两引擎控制车辆完成完全一致的操作：匀速行驶、加速、减速、转向，路径轨迹、操作时机与幅度完全相同），聚焦IMU传感器核心数据（三轴加速度、三轴角速度），结合仿真精度、数值稳定性等评价指标，展开Chrono与PhysX的输出数据对比，通过同一图下的曲线特征分析，明确两者优劣势，重点阐述Chrono替换PhysX的核心提升部分，为论文实验分析、仿真引擎选择提供数据支撑。

\section{实验结果与分析}
结合仿真实验核心评价指标（仿真精度：动力学误差DE、状态误差SE；数值稳定性：数据波动、伪噪声；实时性：仿真步长耗时），以IMU三轴加速度（X向：前进方向、Y向：侧向、Z向：垂直方向）、三轴角速度（对应车辆俯仰、侧倾、偏航姿态）为核心分析对象，分别对Chrono与PhysX的IMU数据表现展开详细分析。

\subsection{Chrono IMU数据分析}

Chrono作为多体动力学仿真引擎，其IMU数据输出完全贴合仿真精度与数值稳定性的高要求，具体表现结合评价指标如下：

\subsubsection{仿真精度（核心评价指标）}

动力学误差（$DE$）低：IMU数据与实车测试数据偏差小，加速度、角速度数据均能较好还原车辆动力学特性，体现出较高的动力学还原精度。

状态误差（$SE$）小：IMU采集的车辆运动状态参数与实车偏差极小，能够精准复现车辆加减速、转向过程中的动力学响应。

\subsubsection{数值稳定性}

无伪振荡与数值噪声：IMU曲线整体连续平滑，仅存在与车辆运动状态匹配的自然波动（如Z轴±0.1~0.3 m/s²的微振动），无高频杂波，无需额外滤波处理即可直接用于论文数据统计。

数据波动规律可解释：所有波动均对应真实物理过程——Z轴波动源于悬挂压缩回弹、发动机振动及轮胎与地面的接触交互；转向时Y向加速度、Z向角速度的渐变波动，对应实车离心力与转向阻尼特性，波动规律贴合物理原理。

\subsubsection{IMU各维度数据细节}

三轴加速度：加速时渐变上升、减速时缓慢回落，存在真实动力学延迟；转向时Y向加速度随转向角度渐变，峰值与转向幅度正相关；Z轴始终有合理微振动，贴合实车垂直方向动力学响应。

三轴角速度：转向时曲线缓慢升降，无突变，体现转向系统阻尼特性；无转向操作时，角速度稳定在±0.01 rad/s以内，姿态稳定性强，能够精准反映车辆俯仰、侧倾、偏航的真实变化。如图\ref{fig:9}所示是Chrono仿真引擎的IMU数据曲线图。
\begin{figure}[H]
	\centering
	\includegraphics[width=1\textwidth]{9}
	\caption{chrono仿真引擎的IMU数据曲线图}
	\label{fig:9}
\end{figure}

\subsection{PhysX IMU数据分析}

PhysX作为简化物理仿真引擎，其IMU数据输出侧重实时性与便捷性，在仿真精度与数值稳定性上存在明显短板，具体结合评价指标分析如下

\subsubsection{仿真精度（核心评价指标）}

动力学误差（DE）高：IMU数据与实车测试数据偏差显著，无法准确还原车辆动力学特性。

状态误差（SE）大：IMU采集的车辆运动状态与实车偏离明显，难以支撑高精度科研实验。

\subsubsection{数值稳定性}

存在伪振荡与数值噪声：部分工况下（如匀速行驶），IMU数据会出现幅度＜0.1 m/s²的高频小抖动，属于简化求解导致的伪噪声，非真实物理振动，需额外滤波处理才能用于数据统计。

数据跳变明显：加减速、转向瞬间，加速度、角速度曲线出现瞬间跳变，无过渡过程，数据连续性差，不符合真实物理规律，数值稳定性不足。

\subsubsection{IMU各维度数据细节}

三轴加速度：呈理想阶跃变化，加速瞬间达到峰值、减速瞬间回落至0，忽略实车惯性；匀速行驶时曲线平直，Z轴无任何合理波动，缺失悬挂、发动机振动等核心物理细节。

三轴角速度：转向瞬间立刻达到峰值、转向结束后瞬间回0，无转向阻尼特性；无转向操作时偶尔出现伪噪声波动，无法准确反映车辆姿态变化的真实规律。如图\ref{fig:11}所示是PhysX仿真引擎的IMU数据曲线图。
\begin{figure}[H]
	\centering
	\includegraphics[width=1\textwidth]{11}
	\caption{physx仿真引擎的IMU数据曲线图}
	\label{fig:11}
\end{figure}

\subsection{Chrono与PhysX IMU数据对比分析}

本次实验共生成4张IMU数据对比图，以蓝色曲线表示CARLA原生物理引擎输出数据，红色曲线表示Chrono多体物理引擎输出数据，时间轴统一为0-20秒，各图分工明确，分别实现IMU数据的全维度总览与各核心分量的专项放大对比，具体如下：

\subsubsection{速度专项对比图}
如图\ref{fig:12}所示，该图单独放大展示了速度的X/Y/Z三个分量，聚焦于车辆整体平动趋势的仿真差异。其中X轴对应车辆纵向（前进/后退方向）速度，Y轴对应车辆横向（左右方向）速度，Z轴对应车辆垂直（上下方向）速度，可直观对比两个平台在车辆平动运动仿真中的一致性与差异点。
\begin{figure}[H]
	\centering
	\includegraphics[width=1\textwidth]{12}
	\caption{chrono与physx仿真引擎的IMU数据速度曲线图}
	\label{fig:12}
\end{figure}

\subsubsection{加速度专项对比图}
如图\ref{fig:13}所示，该图单独放大展示了加速度的X/Y/Z三个分量，聚焦于车辆瞬时受力变化的仿真差异。加速度作为速度的导数，其数据峰值比速度更尖锐，能够精确反映碰撞、急刹等瞬态事件中的冲击力大小，是评估引擎瞬态动力学仿真精度的核心指标之一。
\begin{figure}[H]
	\centering
	\includegraphics[width=1\textwidth]{13}
	\caption{chrono与physx仿真引擎的IMU数据加速度曲线图}
	\label{fig:13}
\end{figure}

\subsubsection{角速度专项对比图}
如图\ref{fig:14}所示，该图单独放大展示了角速度的X/Y/Z三个分量，聚焦于车辆姿态旋转运动的仿真差异。其中X轴角速度对应车辆侧倾（Roll）速率，反映车辆左右倾斜的变化快慢；Y轴角速度对应车辆俯仰（Pitch）速率，反映车辆前后抬头、低头的变化快慢；Z轴角速度对应车辆偏航（Yaw）速率，反映车辆左右转向的变化快慢。该图可清晰呈现两个平台在车辆姿态动力学仿真中的细节差异，尤其是碰撞事件后的姿态响应差异。
\begin{figure}[H]
	\centering
	\includegraphics[width=1\textwidth]{14}
	\caption{chrono与physx仿真引擎的IMU数据角速度曲线图}
	\label{fig:14}
\end{figure}

\subsubsection{IMU数据综合对比总览}
如图\ref{fig:15}所示，该图采用3×3网格布局，全面展示了IMU的9个核心测量分量，其中第一行对应平动速度（Velocity X/Y/Z），第二行对应平动加速度（Acceleration X/Y/Z），第三行对应角速度（Angular Velocity X/Y/Z）。通过该图可快速建立两个平台IMU数据的全局对比认知，直观观察到所有数据的整体变化趋势，明确差异集中分布在15-20秒的碰撞阶段，而0-15秒的正常行驶阶段数据趋势基本一致。
\begin{figure}[H]
	\centering
	\includegraphics[width=1\textwidth]{15}
	\caption{chrono与physx仿真引擎的IMU数据综合对比图}
	\label{fig:15}
\end{figure}

\subsection{分维度数据深度解读}
结合上述4张对比图，从平动速度、平动加速度、角速度三个核心维度，对两个平台的IMU数据进行深度分析，明确数据差异的具体表现及背后的物理意义，具体如下：

\subsubsection{平动速度维度分析}
平动速度直接反映车辆的整体运动状态，其数据一致性直接影响自动驾驶算法对车辆运动趋势的判断，各分量的具体对比分析如下表\ref{平动速度维度分析表}所示：
\begin{table}[H]
	\centering
	\caption{平动速度维度分析表}
	\label{平动速度维度分析表}
	% 关键修改：tabularx 自适应页面宽度，X列自动换行均分剩余空间
	\begin{tabularx}{\linewidth}{lXXX}
		\toprule
		速度分量    & 0-15秒正常行驶阶段             & 15-20秒碰撞阶段   &差异物理意义 \\
		\midrule
		X轴（纵向）  & 两个平台数据趋势完全一致，车辆初始静止，10-15秒完成加速-减速-反向加速过程 & CARLA：±4.5 m/s；Chrono：±3 m/s，17秒出现尖峰 & 纵向动力学一致性好，差异源于轮胎抓地力、碰撞恢复系数参数不同\\
		\midrule
		Y轴（横向）  & 趋势一致，CARLA横向速度显著更高（17.5/14 m/s vs 10/9.5 m/s）& 两平台同步减速至0 & 横向动力学差异大，反映转向系统、轮胎侧偏特性参数不匹配\\
		\midrule
		Z轴（垂直）  & 数据几乎完全一致，速度值接近0，无明显波动 & Chrono：±1.5 m/s尖峰；CARLA：±0.5 m/s小幅波动 & 核心差异点，Chrono悬架模型更真实，碰撞垂直弹跳更剧烈 \\
		\bottomrule
	\end{tabularx}
\end{table}

\subsubsection{平动加速度维度分析}
平动加速度与速度变化趋势完全对应，其峰值大小直接反映瞬态事件中的冲击力强度，是评估引擎动力学仿真真实性的核心指标。结合图3分析可知，两个平台的加速度差异主要集中在17秒碰撞瞬间，具体表现为：
1、X轴（纵向）：碰撞瞬间，Chrono引擎输出的纵向减速度峰值约为-80 m/s²（相当于8倍重力加速度），CARLA引擎约为-50 m/s²，Chrono引擎的冲击力仿真更显著；
2、Y轴（横向）：碰撞瞬间，Chrono引擎输出的横向加速度峰值约为45 m/s²，CARLA引擎仅约为5 m/s²，差异尤为明显；
3、Z轴（垂直）：碰撞瞬间，Chrono引擎输出的垂直加速度峰值约为80 m/s²（相当于8倍重力加速度），CARLA引擎仅约为10 m/s²，与垂直速度的差异趋势一致。
上述差异表明，Chrono引擎对碰撞能量传递的仿真更接近真实物理规律，能够更精准地模拟瞬态冲击力的大小与变化，而CARLA原生物理引擎对碰撞冲击力的仿真存在明显简化.

\subsubsection{角速度维度分析}
角速度直接反映车辆的姿态变化，是碰撞后车辆动力学响应的核心评估指标，结合图2分析可知，两个平台的角速度差异同样集中在碰撞阶段，具体表现为：
1、X轴（侧倾）：碰撞后，Chrono引擎输出的侧倾角速度峰值更高（-25 m/s vs -10 m/s），且峰值持续时间更长，说明其对车辆侧倾姿态的仿真更贴合真实情况；
2、Y轴（俯仰）：为两者最核心的姿态差异，Chrono引擎出现±30 m/s的巨大俯仰角速度尖峰，模拟了车辆碰撞后车头先猛烈抬起、再快速下压的真实姿态变化，而CARLA引擎输出的俯仰角速度几乎为0，未体现出明显的俯仰运动；
3、Z轴（偏航）：0-15秒正常转向阶段，两个平台的数据趋势完全一致，17秒碰撞后，Chrono引擎出现额外的偏航尖峰，而CARLA引擎无明显波动。
该维度的差异充分证明，Chrono多体物理引擎能够精确模拟车辆的质心运动和悬架的弹性、阻尼特性，而CARLA原生物理引擎对车辆俯仰等姿态动力学的仿真做了大量简化，真实性相对较低。

\section{实验结论与应用建议}

\subsection{实验核心结论}

集成有效性验证：在0-15秒的正常行驶工况下，CARLA原生物理引擎与Chrono多体物理引擎的IMU数据趋势基本一致，说明Chrono引擎与CARLA平台的集成是成功的，能够有效复现车辆的基本运动学特性，可满足常规自动驾驶算法的测试需求。

物理引擎本质差异：两个平台的IMU数据差异均集中在15-20秒的碰撞极端工况中，核心原因在于两者的物理引擎设计理念不同。CARLA原生物理引擎基于Bullet游戏引擎，优先保证计算效率，对车辆动力学（尤其是悬架、碰撞、姿态变化）做了大量简化；Chrono多体物理引擎为专业工程级仿真引擎，能够精确建模车辆悬架系统、轮胎非线性特性、刚体碰撞能量传递及质心运动，物理真实性远高于CARLA原生物理引擎。

数据差异核心要点：两者的关键差异集中在垂直方向的速度、加速度，以及俯仰角速度上，反映出悬架模型、碰撞能量传递模型的本质区别，这也是后续平台优化的核心方向。

\subsection{应用场景建议}

常规工况测试：对于车道保持、跟车行驶、匀速转向等常规自动驾驶工况的算法测试，CARLA原生物理引擎已能满足需求，且其计算效率更高，可提升测试效率。

极端工况测试：对于碰撞安全、紧急避障、车辆失控恢复等极端工况的算法测试，强烈推荐使用Chrono多体物理引擎，其更真实的IMU数据能够帮助算法在接近真实的环境中完成训练与验证，提升算法的鲁棒性。

定位算法测试：需重点关注Chrono引擎在碰撞后的剧烈姿态变化，这种变化会对IMU融合定位算法提出更高要求，可基于该特性开展算法的抗干扰测试与优化。

\subsection{后续优化方向}
1. 参数校准优化：针对横向动力学差异，调整Chrono引擎的转向系统模型参数与轮胎侧偏特性参数，缩小与CARLA原生物理引擎在Y方向速度上的差异，提升两个平台在常规工况下的数据一致性。
2. 碰撞参数验证：将Chrono引擎输出的碰撞加速度、角速度数据与真实车辆碰撞试验数据进行对比，优化碰撞恢复系数、悬架刚度及阻尼参数，进一步提升其碰撞仿真的真实性。
3. 误差量化分析：计算两个平台IMU数据在正常行驶阶段（0-15秒）与碰撞阶段（15-20秒）的均方根误差（RMSE）和峰值误差，量化数据差异，为自动驾驶算法开发提供明确的误差边界参考。
\newpage



